\begin{figure}[p]
\centering
\begin{tikzpicture}[
  x=1cm, y=1cm, font=\small,
  node distance=0.5cm,
  box/.style={draw, thick, rectangle, rounded corners=2pt, minimum width=4cm, minimum height=0.8cm, align=center, inner sep=4pt},
  decision/.style={draw, thick, diamond, aspect=2, minimum width=4cm, align=center, inner sep=1pt},
  leaf/.style={draw, thick, rectangle, fill=engblue!10, rounded corners=2pt, minimum width=3.5cm, minimum height=0.8cm, align=center, inner sep=4pt},
  fail/.style={draw, thick, rectangle, fill=engred!10, rounded corners=2pt, minimum width=3.5cm, minimum height=0.8cm, align=center, inner sep=4pt},
  arrow/.style={->, thick, >=stealth}
]

% Start
\node[box] (start) at (0, 14) {Given $\sum a_n$};

% Term test
\node[decision] (term) at (0, 12.5) {$a_n \to 0$?};
\draw[arrow] (start) -- (term);

\node[fail] (diverges1) at (-5, 12.5) {Diverges \\ (term test)};
\draw[arrow] (term) -- node[above, font=\scriptsize] {no} (diverges1);

% Form check: r^n?
\node[decision] (geom) at (0, 11) {form $r^n$?};
\draw[arrow] (term) -- node[right, font=\scriptsize] {yes} (geom);

\node[leaf] (geomtest) at (5, 11) {Geometric: \\ converges iff $|r| < 1$};
\draw[arrow] (geom) -- node[above, font=\scriptsize] {yes} (geomtest);

% Factorial?
\node[decision] (fact) at (0, 9.5) {factorial?};
\draw[arrow] (geom) -- node[right, font=\scriptsize] {no} (fact);

\node[leaf] (ratio) at (5, 9.5) {Ratio test};
\draw[arrow] (fact) -- node[above, font=\scriptsize] {yes} (ratio);

% Alternates?
\node[decision] (alt) at (0, 8) {alternates?};
\draw[arrow] (fact) -- node[right, font=\scriptsize] {no} (alt);

\node[leaf] (altseries) at (5, 8) {Alternating-series test \\ ($b_n \downarrow 0$ gives convergence)};
\draw[arrow] (alt) -- node[above, font=\scriptsize] {yes} (altseries);

% 1/n^p like?
\node[decision] (poly) at (0, 6.5) {$\sim 1/n^p$ large $n$?};
\draw[arrow] (alt) -- node[right, font=\scriptsize] {no} (poly);

\node[leaf] (polytest) at (5, 6.5) {(Limit) comparison \\ with $\sum 1/n^p$};
\draw[arrow] (poly) -- node[above, font=\scriptsize] {yes} (polytest);

% Integrable f(n)?
\node[decision] (intf) at (0, 5) {clean $\int f$ exists?};
\draw[arrow] (poly) -- node[right, font=\scriptsize] {no} (intf);

\node[leaf] (inttest) at (5, 5) {Integral test};
\draw[arrow] (intf) -- node[above, font=\scriptsize] {yes} (inttest);

% Root test
\node[decision] (rootq) at (0, 3.5) {$n$-th-root limit?};
\draw[arrow] (intf) -- node[right, font=\scriptsize] {no} (rootq);

\node[leaf] (rootest) at (5, 3.5) {Root test};
\draw[arrow] (rootq) -- node[above, font=\scriptsize] {yes} (rootest);

% Inconclusive
\node[fail] (incon) at (0, 2) {Inconclusive: \\ split into pieces, \\ or look up};
\draw[arrow] (rootq) -- node[right, font=\scriptsize] {no} (incon);

\end{tikzpicture}
\caption[Convergence-test decision tree]{The convergence-test ladder
as a decision tree. The engineer descends from cheapest test to most
specialised. The term test catches the simplest divergent series; the
geometric form is recognised on sight; the ratio test handles factorials;
the alternating-series test handles oscillating terms with monotone
magnitudes; comparison handles polynomial-tail series; the integral test
handles cases where $f(n)$ has a clean antiderivative; the root test is
the back-stop. The leaves at the right are the tests; the diamonds are
the recognition conditions. Engineers run the tree in order until a leaf
returns a verdict.}
\label{fig:vol02:ch04:test-decision-tree}
\end{figure}
